\documentclass[../../../include/open-logic-section]{subfiles}

\begin{document}

\olfileid{sth}{choice}{woproblem}
\olsection{સુક્રમિતતાની સમસ્યા}

સુક્રમિતતા સ્વીકારીએ કે નહીં તેના પર ઘણું
નિર્ભર છે, એ સ્પષ્ટ છે. પરંતુ આ સિદ્ધાંતની
ચર્ચા સામાન્ય રીતે તેના સમકક્ષ બીજા સિદ્ધાંત,
પસંદગીની સ્વયંસિદ્ધિ, પર કેન્દ્રિત રહી છે.
તેથી હવે તેના તરફ વળીએ (અને સમકક્ષતા
સાબિત કરીએ).

\citeyear{Cantor1883}માં કૅન્ટરે સુક્રમિતતાની
સ્વયંસિદ્ધિને સમર્થન આપતાં તેને ``વિચારનો
એક એવો નિયમ કહ્યો જે મને મૂળભૂત,
પરિણામોની દૃષ્ટિએ સમૃદ્ધ અને ખાસ કરીને
સર્વસાધારણ માન્યતાને કારણે નોંધપાત્ર લાગે છે''
(\citeauthor{Potter2004}ના
\citeyear[p.~243]{Potter2004}માં અવતરણ).
પરંતુ અંતે કૅન્ટરને ખાતરી થઈ કે આ
``સ્વયંસિદ્ધિ''ને પુરાવાની જરૂર છે.
ગણિતજ્ઞ સમુદાય પણ એમ જ માનવા લાગ્યો.

\citeyear{Zermelo1904}માં ઝર્મેલોએ આ
સમસ્યા ``ઉકેલી''. તેમનો ઉકેલ સમજાવવા
થોડી વ્યાખ્યાઓ જોઈએ.

\begin{defn}
વિધેય $f$ \emph{પસંદગી વિધેય} ત્યારે અને
ત્યારે જ છે, જ્યારે દરેક $x \in \dom{f}$
માટે $f(x) \in x$. આપણે $f$ને
\emph{$A$ માટેનું પસંદગી વિધેય} ત્યારે કહીએ,
જ્યારે $f$ પસંદગી વિધેય હોય અને
$\dom{f} = A \setminus \{\emptyset\}$.
\end{defn}

સહજ રીતે, દરેક અરિક્ત ગણ $x \in A$ માટે
$A$ માટેનું પસંદગી વિધેય $x$માંથી એક ચોક્કસ ઘટક
$f(x)$ \emph{પસંદ કરે છે}. હવે પસંદગીની
સ્વયંસિદ્ધિ આ છે:

\begin{axiom}[પસંદગી]
	દરેક ગણ માટે પસંદગી વિધેય છે.
\end{axiom}

ઝર્મેલોએ બતાવ્યું કે પસંદગીમાંથી સુક્રમિતતા
મળે છે, અને ઊલટું પણ:

\begin{thm}[$\ZF$માં]\ollabel{thmwochoice}
સુક્રમિતતા અને પસંદગી સમકક્ષ છે.
\end{thm}

\begin{proof}
\emph{ડાબેથી જમણે.} $A$ને ગણોનો ગણ લો.
યોગગણની સ્વયંસિદ્ધિથી $\bigcup A$ છે;
તેથી સુક્રમિતતાથી $\bigcup A$નો કોઈ
સુક્રમ $<$ છે. હવે $f(x) = \text{$x$નો
$<$ મુજબ લઘુતમ ઘટક}$ લો. આ $A$ માટેનું
પસંદગી વિધેય છે.

\emph{જમણેથી ડાબે.} $A$ નિશ્ચિત લો.
ખાલી ગણ પહેલેથી સુક્રમિત છે; તેથી હવે
ગણ અરિક્ત છે એમ ધારો.
પસંદગીથી $\Pow{A} \setminus
\{\emptyset\}$ માટે પસંદગી વિધેય $f$
છે. અનંતાતીત પુનરાવર્તનથી એક વિધેય
આ રીતે વ્યાખ્યાયિત કરો:
\begin{align*}
	g(0) &= f(A)\\
	g(\alpha) &= 
		\begin{cases}
			\text{થોભો!{}} &\text{જો }A = \funimage{g}{\alpha}\\
			f(A \setminus \funimage{g}{\alpha}) & \text{અન્યથા}\\	
		\end{cases}
\end{align*}
``થોભો!'' એ વધુ લાંબી વ્યાખ્યાને બદલે
વપરાયેલ ટૂંકું સૂચન છે. એટલે પહેલી વાર
$A = \funimage{g}{\alpha}$ થાય ત્યારે,
બધા $\delta \geq \alpha$ માટે
$g(\delta)=A$ મૂકો. શરૂઆતમાં આથી
ફક્ત એક \emph{પદ} વ્યાખ્યાયિત થાય છે
એટલું જ નિશ્ચિત છે
(\olref[sth][spine][recursion]{transrecursionschema}
પછીની ટિપ્પણીઓ જુઓ); પરંતુ આપણે
બતાવીશું કે વિરામ પહેલાં તે ખરેખર વિધેય છે.
\sourcecorrection{OLMD-147}{મૂળમાં
$g(0)=f(A)$ પહેલાં $A$ અરિક્ત હોવાનું
કહ્યું નથી. $A=\emptyset$ હોય તો
$f$નું ક્ષેત્ર તેને સમાવતા નથી; તે
કિસ્સો સીધો ઉકેલીને પુનરાવર્તન
અરિક્ત $A$ માટે જ શરૂ કર્યું છે.}

$f$ પસંદગી વિધેય હોવાથી, વિરામ પહેલાંના
દરેક $\alpha$ માટે
$g(\alpha) = f(A \setminus
\funimage{g}{\alpha}) \in A \setminus
\funimage{g}{\alpha}$, એટલે
$g(\alpha) \notin \funimage{g}{\alpha}$.
તેથી વિરામ પહેલાં
$g(\alpha) = g(\beta)$ હોય, તો
$g(\beta) \notin \funimage{g}{\alpha}$,
એટલે $\beta \notin \alpha$; તેવી જ રીતે
$\alpha \notin \beta$. ત્રિવિભાજનથી
$\alpha = \beta$. આમ વિરામ પહેલાંનું
$g$ !!{injective} છે.

હવે નોંધો કે ``થોભો!'' ખરેખર થાય છે;
એટલે એવી કોઈ લઘુતમ ક્રમવાચક સંખ્યા
$\alpha$ છે કે $A = g[\alpha]$.
ઊલટું ધારો; તો દરેક ક્રમવાચક સંખ્યા
$\alpha$ માટે વિરામ પહેલાંનું $g$
!!{injective} હોવાથી
$\cardless{\alpha}{\Pow{A} \setminus
\{\emptyset\}}$ મળે, જે
\olref[hartogs]{HartogsLemma}નો વિરોધ
કરે છે. તેથી
$\funimage{g}{\alpha}=A$ પણ છે.

આ બધું ભેગું કરતાં
$\funrestrictionto{g}{\alpha}$ એ કોઈ
ક્રમવાચક સંખ્યામાંથી $A$ પર
!!a{bijection} છે. હવે આ $g$ના વિરામ
પહેલાંના ભાગથી
$A$ને સુક્રમિત કરી શકાય.
\sourcecorrection{OLMD-148}{મૂળમાં વિરામની
પહેલી ઘડીએ બધા $\delta\leq\alpha$ માટે
$g(\delta)=A$ લખાયું છે; એમ કરવાથી
અગાઉની પસંદગીઓ જ નષ્ટ થાય. અહીં
$\delta\geq\alpha$ કર્યું છે અને
અંતઃક્ષેપ/એક-એક વ્યાપ્તતાનો દાવો ફક્ત
વિરામ પહેલાંના $\funrestrictionto{g}{\alpha}$
માટે કર્યો છે; વિરામ પછીનું સ્થિર મૂલ્ય
આ દાવામાં આવતું નથી.}
\end{proof}

આથી સુક્રમિતતા અને પસંદગી બંને સાથે
ઊભાં રહે છે અથવા સાથે પડે છે. પણ પ્રશ્ન
હજી છે: તેઓ ઊભાં રહે છે કે પડે છે?

\end{document}
